\section{Model scope}\label{sec:scope}
NetSy synthesizes freight-related network loads and emissions from socioeconomic
and spatial information. Its scope connects the activities that generate demand
for goods with the vehicle movements needed to serve that demand and the effects
of those movements on the road network and the environment.

Seven functional responsibilities describe this scope without prescribing a
particular algorithm or software module:
\begin{enumerate}
\item \textbf{Study-area representation:} characterize economic actors,
locations, facilities, vehicles and the spatial and network context.
\item \textbf{Goods demand:} represent the quantities produced, consumed or
requested by those actors.
\item \textbf{Goods exchanges:} establish the origins, destinations and quantities
of exchanges between actors or areas.
\item \textbf{Logistics organization:} represent how exchanges become shipments,
distribution chains and transport services.
\item \textbf{Vehicle activity:} organize those services into movements and tours,
including timing, capacity use and empty travel.
\item \textbf{Network loading:} locate vehicle movements on roads under specified
route-choice and traffic conditions.
\item \textbf{Environmental assessment:} relate vehicle use and operating
conditions to energy consumption and emissions.
\end{enumerate}
These responsibilities provide a common structure for discussing alternative
formulations and study areas. Their boundaries overlap: establishment attributes
inform demand and logistics choices, for example, while vehicle activity informs
both network loading and emissions. Data feasibility, calibration, uncertainty
and validation are considered throughout. The scope describes what NetSy is
intended to address; the review does not select or evaluate its implemented
algorithms.

\section{Literature review}\label{sec:slr}
The review establishes the methodological and data foundations for these
functional responsibilities. Evidence identification uses eight complementary
search strands; the findings are synthesized by the seven responsibilities in
Section~\ref{sec:scope}. Search organization and synthesis organization therefore
serve different purposes. An established foundation means a recognized modelling
principle, not a universally superior algorithm.

MASS-GT, from which NetSy originated, provides a designated reference
\parencite{massgt2018,massgt2025}. Its literature is identified separately and
interpreted within each functional responsibility, alongside the wider academic
evidence and the information needed to apply it. This does not treat
reproduction of MASS-GT as evidence of external validity. Positioning NetSy's
implementation and numerical results is deferred to subsequent work.

\subsection{Methodology}\label{sec:slr-methodology}
\paragraph{Questions and scope.}
Three questions guide the synthesis:
\begin{enumerate}
\item Which approaches address the seven functional responsibilities, and what theoretical,
comparative and validation evidence supports them?
\item Which data configurations, classifications and resolutions support these
approaches across study areas, and what adaptations are possible when
information is missing?
\item Which formulations, inputs and assessment evidence define the MASS-GT
reference within each functional responsibility?
\end{enumerate}
The scope includes urban, metropolitan and regional freight applications,
with economic, network and environmental foundations where they inform those
applications. The main search language is English; Swedish official data
documentation is also examined. A versioned protocol governs the present
reassessment. It was written after exploratory retrieval, so the earlier
searches are not represented as prospectively protocol-controlled. The protocol
and extraction approach draw on \textcite{kitchenham2007}; reporting is organized
with reference to PRISMA-ScR and PRISMA-S
\parencite{prismascr2018,prismas2021}. These references describe guidance used,
not a claim that every recommended safeguard has been completed.

\paragraph{Search strands and synthesis framework.}\label{sec:review-framework}
Six subject strands address demand generation and economic representation;
spatial exchange and origin--destination (OD) inference; logistics decisions and
shipment formation; vehicle activity, tours and scheduling; network models and
route choice; and energy and emissions. A seventh strand examines data
availability, representation and transferability. An eighth identifies MASS-GT
publications and documentation. These are retrieval strands, not eight equivalent
academic fields. Data evidence and the reference-model literature inform the
functional synthesis throughout; they do not form separate result sections.

Table~\ref{tab:review-framework} records the principal connections. The framework
is many-to-many: a source can inform several responsibilities regardless of
which query retrieved it. The functional descriptions organize interpretation;
they are not inclusion rules requiring agreement with NetSy or MASS-GT.
\begin{table}[htbp]
\centering\small
\caption{Search strands and their principal contribution to the functional synthesis.}
\label{tab:review-framework}
\begin{tabularx}{\linewidth}{@{}p{18mm}p{65mm}Y@{}}
\toprule Strand & Search focus & Principal synthesis sections \\
\midrule
S1 (A1) & Demand generation and economic representation & Study-area representation; goods demand \\
S2 (A2) & Spatial exchange and OD inference & Goods demand; goods exchanges \\
S3 (A3) & Logistics decisions and shipment formation & Goods exchanges; logistics organization \\
S4 (A4) & Vehicle activity, tours and scheduling & Logistics organization; vehicle activity \\
S5 (A5) & Network models and route choice & Network loading; vehicle activity \\
S6 (A6) & Energy and emissions & Environmental assessment; network loading \\
S7 (A7) & Data availability, representation and transferability & All seven responsibilities \\
S8 (M) & MASS-GT reference literature and documentation & All seven responsibilities \\
\bottomrule
\end{tabularx}
\end{table}
A1--A7 and M are the historical query identifiers retained in the appendix and
metadata. S1--S8 identify the eight strands in the present framework. This
mapping preserves the retrieval history; it does not rename old searches or
imply they were executed under a prospectively defined functional framework.

Framework-based narrative synthesis groups the extracted findings by functional
responsibility. Each section considers available approaches and assessment
evidence, their data requirements and access conditions, and the corresponding
MASS-GT configuration. The same source may support several sections without
becoming several independent studies. Appendix~\ref{app:search} retains the
literal searches and selection record; a supporting source-to-function index
links the cited evidence to the synthesis sections.

\paragraph{Identification and selection.}
OpenAlex and Crossref retrieval took place on 11, 12 and 16 November 2026.
The first two rounds returned at most 25 ranked records per engine and query;
the third expanded the window to 100. No lower publication-year limit was used.
The third round added the data-availability search strand and separate MASS-GT discovery. Official
bibliographies, targeted reference tracing, institutional repositories and
technical guidance supplemented the database searches. Engine-specific strings,
dates, limits and supplementary routes are reproduced in
Appendix~\ref{app:search}. OpenAlex search and Crossref bibliographic ranking
are different retrieval mechanisms; equivalent Boolean behaviour is not assumed.

The initial available-document cohort comprised 52 PDFs, including two one-page
partial files. Institutional retrieval on 18 November supplied twelve files:
eleven complete relevant papers and an unrelated course syllabus, excluded at
eligibility checking. Two complete papers replace the earlier partial files.
The current cohort therefore contains 61 complete source documents: 59 retained
as substantive evidence and two as background. Sources are eligible when they provide
an identifiable formulation, assessment or data specification relevant to a
functional responsibility or its information requirements. Peer-reviewed studies, author versions, theses, books and
official guidance are distinguished, as are conceptual proposals and empirical
findings. Generic shortest-path exposition and a discovery compendium are
background rather than substantive findings.

Selection was originally purposive within bounded ranked results; the remaining
database records have not all undergone eligibility screening. The reassessment
applies explicit criteria to the held-document cohort rather than claiming to
have screened every search hit. Twenty-eight access requests remain unresolved,
including three of the fourteen institutional priorities. They are not excluded
on scientific grounds, and abstract-only sources do not support methodological
performance claims. Appendix~\ref{app:selection} reports the separate retrieval
and document counts without constructing a fictitious completed selection flow.
The evidence cohort is unchanged from release v0.1.0. The present revision
reorganizes its synthesis; it introduces neither new searches nor new selection
decisions.

\paragraph{Extraction and appraisal.}
For each retained source, extraction records the outcome and formulation, actors
and boundaries, inputs, classifications and resolution, access conditions,
estimation or calibration, validation design, findings and limits of inference.
Targeted relevant sections are inspected and their PDF pages recorded. Matching
calibration controls, predicting held-out observations, transferring between
areas and validating against independent measurements are treated as different
evidence. Optimization performance is distinguished from behavioural prediction;
model intercomparison from measured accuracy. Official data specifications
establish definitions and access arrangements, not model performance.

Data appraisal distinguishes inputs needed to run a formulation, observations
used to estimate or calibrate it, and independent validation information.
Meaning, units, classification, spatial and temporal resolution, coverage,
vintage, access restrictions and transfer assumptions are examined within each
responsibility. A nominally available dataset may not support the required
resolution or population.

Narrative synthesis separates established principles, alternative formulations
and unresolved questions. Heterogeneous outcomes are not pooled into a common
accuracy score, and related publications or shared datasets are not counted as
independent votes. Retrieval, reassessment and writing are AI-assisted, without
independent duplicate screening or extraction. Bounded coverage, uneven access,
English-led searching and the absence of independent checking limit the
conclusions. This is an available-evidence scoping synthesis, not an exhaustive
review or a measured consensus of the entire field.

\paragraph{Reference configuration.}
The principal published anchor is the full-system account of
\textcite{massgt2025}. Earlier publications document different stages, including
an initial prototype and a zero-emission-zone application
\parencite{massgt2018,massgt2022}. Technical documentation version~3.0
\parencite{massgtmanual2023} is retained at a specified code checkout. The
published configurations are not assumed to coincide exactly with that commit.
The reference is interpreted within each functional responsibility. Its access
conditions are examined alongside those of the wider literature, rather than
being assumed to apply to all study areas.

\subsection{Findings by functional responsibility}\label{sec:slr-synthesis}

\subsubsection{Study-area representation}\label{sec:f1}
Study-area representation defines the actors, places, facilities, fleets and
spatial units to which the remaining formulations refer. The literature connects
these choices to the intended outcome and available observations. Aggregate
commodity models, regional logistics models and urban touring models need
not represent the same actors or spatial detail \parencite{chow2010}. The FHWA
source inventory distinguishes industry and commodity classifications, spatial
and temporal detail, and the roles of observations in modelling
\parencite{fhwa2018}. These documented US configurations are informative
examples, not a worldwide catalogue of currently available inputs.

When individual firms must be represented but are not directly observed,
synthetic populations provide an additional approach. Statistical reconstruction
uses aggregate information and available samples to construct establishments
with relevant attributes, as illustrated by \textcite{wisetjindawat2009}.
The available evidence does not establish a universal reconstruction winner.
Firm-specific head-to-head comparisons remain an access gap, and a synthetic
population matching its construction totals does not independently validate
firm-level freight behaviour.

\paragraph{Data requirements and feasibility.}
An establishment register, sectoral employment totals, household statistics,
facility locations and fleet records describe different populations and
attributes. The required representation depends on which of them can be
observed and which must be synthesized. National aggregates can support broad
analysis without identifying local establishments or their exact locations
\parencite{fhwa2017}.

Data fusion can enrich the representation without making all attributes
observed. The Dutch study of \textcite{mohammed2023} classifies logistics nodes
using business registries, postal geography and make/use information, with
certainty indicators for inferred attributes. Its XML sample favours larger,
technologically equipped carriers; single-employee units are excluded from
its firm classification. More detail therefore does not remove selection bias
or establish access to the underlying records. The established requirement is
an explicit account of the represented population and the support for each
attribute, rather than maximum synthetic detail.

\paragraph{MASS-GT reference.}
The published Dutch configuration uses zonal socioeconomic information and
logistics locations to represent firms, households and distribution nodes,
including employment and supply--use relationships \parencite{massgt2025}.
These inputs support its particular actor and spatial configuration. The paper
aims to use standard public inputs where possible, but states that the authors
lack permission to share the application data. The reference is consequently
a mixed-access representation, not evidence that equally detailed actor,
facility and fleet information is available elsewhere.


\subsubsection{Goods demand}\label{sec:f2}
A central distinction in this literature is between the amount of goods produced
or consumed and the vehicle visits associated with moving those goods. Freight
trip generation depends on logistics decisions as well as economic activity:
consolidation and larger shipments can accommodate more goods without a
proportional increase in trips. The NCFRP guidebook therefore treats goods and
trip generation as separate concepts and emphasizes differences between economic
activities \parencite{ncfrp2012}. A single trip rate applied uniformly across
businesses can conceal those differences.

Establishment-level statistical models commonly relate freight activity to
industry, employment, floor area and other business characteristics. The
literature supports choosing these relationships according to the outcome and
the observations, rather than assuming one universal specification.
\Textcite{oliveira2022} compare delivery-generation models across Brazilian
municipalities, while \textcite{pani2021} examine geographical transfer and local
updating in India. The latter find that transferability varies by sector, region
and whether goods quantities or trips are being estimated. The broader review of
\textcite{mcleod2019} places generation within changing business strategies,
land-use patterns and policy contexts, which can affect the usefulness of an
otherwise satisfactory statistical relationship.

The manufacturing-establishment study of \textcite{balla2023} adds evidence on
transfer direction and local updating. It uses surveys from four Indian regions
and geographical covariates to propose reduced sampling in similar areas.
This supports examining data-saving transfer strategies, but the manufacturing
sample and extrapolation to unsurveyed cities delimit the claim. A proposed
sample fraction should not become a universal input requirement.

The established foundation is to distinguish economic demand from vehicle
activity, account for differences between establishments and sectors, and
validate generation estimates at the scale where they will be used.

\paragraph{Data requirements and feasibility.}
Generation relationships need observations that match their outcome: goods
quantities, deliveries and vehicle visits are not interchangeable dependent
variables. Public survey descriptions also differ from accessible microdata.
Eurostat's road freight manual specifies vehicle, journey and goods information,
but collection, confidentiality and dissemination are separate processes
\parencite{eurostat2025}. Sweden's commodity-flow surveys combine selected-sector
sampling with administrative sources; establishment stratification uses industry
and employment size \parencite{trafa2022,trafa2026method}. The 2025 quality
declaration limits interpretation of detailed geography through coverage and
confidentiality, and research access to microdata requires assessment
\parencite{trafa2026quality}. A published survey does not imply that a local
model can obtain precise establishment-level generation observations.


\paragraph{MASS-GT reference.}
Conventional goods demand in the published framework is conditioned on an
aggregate commodity matrix, disaggregated to shipments and economic actors;
parcel demand is represented separately from households and businesses
\parencite{massgt2018,massgt2025}. Zonal employment and supply--use information
support the actor representation. This distinction matters: synthesizing
shipments consistent with aggregate demand is a different task from independently
forecasting that demand from socioeconomic attributes alone.

\subsubsection{Goods exchanges}\label{sec:f3}
Input--output approaches represent economic relationships between sectors and
regions. They are useful when the research question concerns production,
consumption or changes in economic structure. Their use in freight modelling
requires additional information about spatial trade and the conversion of
monetary values into physical goods. \Textcite{li2023} illustrate this connection
in an interregional application in China. Its interpretation remains conditional
on sector coverage, economic assumptions and conversion data.

Spatial-interaction approaches allocate flows according to the activity at
origins and destinations and the difficulty or attractiveness of exchanging
between them. Gravity and entropy-based approaches are established references
when production and attraction totals are available \parencite{sheffi1985}.
They still require evidence about how distance, cost and competing destinations
affect exchange. \Textcite{celik2001} examine these influences in commodity
flows. The interdisciplinary review of \textcite{rong2024} also distinguishes
spatial-interaction approaches from the construction, estimation and forecasting
of OD movements. Its largely human-mobility evidence, however, cannot be assumed
to apply directly to freight.

Estimating OD movements from traffic counts, trajectories or other partial
observations presents a different challenge. Several movement patterns may
explain the same observations, and the result depends on assumptions about
routes, vehicle use and data coverage. \Textcite{raothanachonkun2007} demonstrate
the importance of trip chains and empty movements when translating commodity
flows into truck trips. In a Chilean highway application, \textcite{basso2022}
combine multiple data sources to infer truck origins and destinations, but the
coverage and destination assumptions remain tied to the available samples.

\Textcite{holguin2008} connect a multicommodity loaded-demand formulation to
an empty-trip model and assess both against survey-based OD and traffic
information from Guatemala. Explicit empty movements improve count replication
in the tested formulations. However, the gravity parameters can remain poorly
recovered even when network counts fit better. The comparison therefore supports
representing empty movements while also showing why a count fit cannot uniquely
identify the underlying commodity relationships.

Classification and calibration also affect supplier-level models.
\Textcite{ismael2026}, an author preprint, combine aggregate commodity controls,
shipping-distance information and industry--commodity mappings across four US
metropolitan areas. Their treatment makes the information lost through commodity
aggregation explicit. Some establishment-synthesis inputs can be proprietary,
and matching the flows used as controls is not independent validation of actual
supplier relationships.

The literature therefore supports distinguishing economic exchange, spatial
allocation and inference from incomplete observations. Balancing known totals is
a useful starting point when those totals are credible; it does not establish
that every OD flow is correct. Validation must also consider flow patterns,
distances, empty vehicles and the representativeness of the observations.

\paragraph{Data requirements and feasibility.}
OD inference can use commodity surveys, trade matrices, counts and trajectories,
but these observe different objects. GPS traces can reveal movements while
omitting commodities, quantities and purposes; goods surveys may not identify
complete vehicle sequences \parencite{shoman2023}. Classification crosswalks
between industry, commodity and vehicle categories can be many-to-many and
require aggregation or assumptions about mixed production. The supplier-selection
preprint of \textcite{ismael2026} illustrates a US industry--commodity mapping
and a deliberate reduction of commodity detail. Synthetic bilateral detail
should remain distinguishable from observed exchange information.

The European corridor application of \textcite{petri2016} offers an aggregate-data
configuration with a different forecasting target from urban tour
reconstruction. Its feasibility cannot be interpreted as evidence that the same
inputs identify fine-grained urban supplier relationships.


\paragraph{MASS-GT reference.}
The 2025 configuration allocates commodity flows over distribution channels and
assigns shipments to firms or logistics nodes. Industry make/use probabilities,
actor size and transport costs influence this disaggregation
\parencite{massgt2025}. The aggregate regional matrix is supplied by BasGoed in
the Dutch application. MASS-GT therefore anchors a controlled disaggregation
configuration; it does not establish that equivalent OD controls can be obtained
or inferred with equal reliability in another study area.

\subsubsection{Logistics organization}\label{sec:f4}
Logistics research treats shipment size, delivery frequency, inventory and
transport service as connected decisions. Larger shipments may reduce the
frequency and cost of transport while increasing storage requirements or
reducing service flexibility. The review of \textcite{tavasszy2020} organizes
these choices by actors and functional areas, emphasizing dependencies along the
supply chain. Inventory theory provides a widely used explanation of the
trade-off between replenishment, storage and transport, developed empirically
and in relation to distribution networks by \textcite{combes2012,combes2016}.

One established approach assumes that actors select the arrangement with the
lowest logistics cost, subject to service and operational requirements. Another
estimates observed choice behaviour, allowing differences between firms and
factors not captured by measured costs. \Textcite{abate2019} compare stochastic
and deterministic logistics modelling in the Swedish national freight context.
Their different responses to cost changes show why a base-year match alone does
not settle how a model will respond to an intervention.

The literature does not establish that all connected decisions must be represented
by the same statistical structure. \Textcite{ahmed2023} compare joint and
sequential shipment-size and vehicle choices; a shipment-size-first sequential
structure fits their dataset best, while their results also indicate dependence
between the choices. \Textcite{piendl2022} investigate segments defined by
logistics requirements rather than administrative commodity categories. The
prediction gains are modest overall and differ between shipment classes.
\Textcite{odolinski2024} further show how missing logistics information and the
representation of costs affect estimation.

The established position is that shipment formation should reflect the
interdependence of logistics choices and the circumstances of decision makers.
Cost-based, behavioural, joint and sequential approaches are alternatives whose
suitability depends on data and purpose. Stronger evidence requires checks of
shipment sizes, frequencies, service choices and responses to changing
conditions, rather than selecting a method solely from its fit to one dataset.

\paragraph{Data requirements and feasibility.}
Shipment records, business attributes and logistics-node information support
estimation only to the extent that the relevant alternatives and costs are
observed or credibly reconstructed. \Textcite{odolinski2024} show why missing
logistics information and cost representation matter. Transferred coefficients,
assumed channel shares and estimated choice parameters should therefore be
reported separately. Fit under existing costs does not establish how actors
will adopt a new hub, vehicle technology or collaborative service.


\paragraph{MASS-GT reference.}
The 2025 specification combines an empirically derived distribution-channel
split with a simultaneous shipment-size/vehicle-type choice model motivated by
transport and inventory costs, followed by delivery-time choice
\parencite{massgt2025}. The earlier zero-emission application used a less developed
choice configuration \parencite{massgt2022}. These stages must remain distinct.
Estimated logistics choices and expert-defined responses to new services also
provide different kinds of evidence; a scenario response is not automatically
an observed adoption relationship.

The empirical interpretation of this reference depends on the attributes
available for estimation. \Textcite{mohammed2023} enrich Dutch shipment diaries
with firm and logistics-node information and explicitly distinguish deterministic
from probabilistic classifications. These inferred logistics roles support
analysis of distribution chains, but are not equivalent to directly observing
every supplier, receiver or decision maker.

\subsubsection{Vehicle activity}\label{sec:f5}
Two broad purposes shape this literature. Descriptive models seek to reproduce
how vehicles are used; optimization models seek efficient plans for serving a
given demand. Vehicle-routing research provides established approaches to
planning deliveries under capacity, timing and fleet restrictions.
\Textcite{vidal2013} review how these requirements shape routing problems and
the methods used to solve them. A routing heuristic is a way of finding a plan;
it is not, by itself, an explanation of observed delivery behaviour.

Descriptive tour models represent decisions such as grouping shipments, adding
stops and ending a tour. \Textcite{thoen2020} estimate shipment-based tour
formation with operational constraints and assess its reproduction of observed
tour characteristics. Their definition excludes empty movements, illustrating
why tour definitions must be reconciled before results are compared. A model
can reproduce recorded delivery tours while leaving other vehicle activity
unexplained.

Detailed operational simulation provides another perspective. In the
pickup-and-delivery study of \textcite{lyu2021}, GPS and operator information
support comparisons between network and process representations. Additional
detail improves the reproduction of aspects of that case, but a known service
sequence is different from predicting the sequence of a new operator. Aggregate
trip-chain models and inverse optimization are further approaches, particularly
where observed activity must be reconciled with operational constraints.

A further approach estimates the parameters of an optimization-based tour model
from observed activity. \Textcite{nadi2024} use Bayesian optimization with
partially observed Dutch carrier tours to reproduce spatial and temporal
properties. This makes optimization a calibrated descriptive tool, rather than
assuming that a shortest plan explains behaviour. Its required carrier records
are permission-controlled, and the authors distinguish reproduced patterns from
knowledge of the causal reasons behind planners' choices.

A different information regime is demonstrated by \textcite{davis2024}.
Their New York framework generates candidate truck tours and balances flows to
aggregate controls, using transferred generation relationships and local network
and count information. It can represent truck agents without reconstructing
individual shipper relationships. The assessment is at boundary-flow and temporal
scales; parcel vans are omitted, service times are assumed, and some constants
are locally adjusted. It consequently supports a mesoscopic alternative, not an
interchangeable reconstruction of observed shipment-level operations.

Transfer of a descriptive framework also requires local work.
\Textcite{ferguson2012} adapt Calgary submodels to Greater Toronto and Hamilton,
combining local commercial firm data, auxiliary movements and recalibrated
constants. Their traffic comparisons support a particular transferred
configuration; they do not demonstrate universal coefficient portability or
accurate reproduction of each individual tour.

Across these approaches, the recurring requirements are feasible service,
consistent vehicle loads, appropriate timing and explicit treatment of empty
movements. Predictive models are assessed against observed activity, whereas
planning methods are also assessed by feasibility, cost and computational
performance. A shorter optimized route does not automatically provide a better
prediction of what carriers will do. The suitable approach depends on which of
these questions the study addresses.

\paragraph{Data requirements and feasibility.}
Detailed operational models can require carrier diaries, GPS, service times,
vehicle capacities and activity purposes. Permission-controlled records support
the preference-based tour model of \textcite{nadi2024}; they are not unrestricted
inputs for prospective users. The development proposal of \textcite{moore2017}
also distinguishes planned GPS access from confidential shipment information:
a proposed collection is not evidence of a completed, validated model.

Three adaptation strategies are visible in the inspected cases: borrow empirical
components and calibrate locally, as in SimMobility Freight
\parencite{sakai2020}; transfer a tour framework with local re-estimation
\parencite{ferguson2012}; or change the representation to use aggregate controls
\parencite{davis2024}. The latter's public-data objective is qualified by old
input vintages, assumed service parameters, adjusted coefficients and
study-specific tract aggregations. These strategies support different questions;
they are not interchangeable reconstructions of individual operations.


\paragraph{MASS-GT reference.}
The published shipment scheduler combines estimated tour-ending decisions with
feasible shipment selection, proximity-based insertion and local improvement
of stop order. Parcel scheduling instead groups delivery zones under vehicle
capacity and improves the stop sequence \parencite{massgt2025}. This is a hybrid
of behavioural estimation and heuristics. The tour-formation evidence of
\textcite{thoen2020} supports a particular activity definition, so the coverage
of empty and service movements still needs reconciliation when assessing a
complete vehicle inventory. The versioned manual additionally documents service
trip simulation \parencite{massgtmanual2023}.

Microhub operations are part of the published MASS-GT/TFS lineage.
\Textcite{microhubs2024} model feeder and last-yard tours under nine Rotterdam
configurations combining hub locations, vehicle types and carrier collaboration.
The HARMONY synthesis places these alongside zero-emission zones, crowdshipping
and land-use scenarios \parencite{harmony2024}. Their results show how assumptions
about organization alter operations, but collaboration and adoption are imposed
scenario conditions rather than estimated responses. These publications are
related applications, not independent replications of the same policy effect.

\subsubsection{Network loading}\label{sec:f6}
Network-assignment research explains how a given set of vehicle movements is
distributed across roads. The simplest established reference assigns movements
to the least-cost paths at supplied travel times and costs. Other approaches
represent variation in route preferences or allow congestion to change travel
costs as vehicles enter the network. Equilibrium methods describe route choices
under those interactions, while dynamic approaches also represent how traffic
conditions evolve over time \parencite{sheffi1985,wang2018}.

The appropriate level of detail depends on the question. Supplied speeds may be
adequate for estimating loads under specified operating conditions. An
intervention intended to change congestion requires attention to the surrounding
traffic and how travel conditions respond. \Textcite{hwang2021} illustrate a
freight-assignment application that incorporates background traffic. Where
queues, departure timing or spillback are material, dynamic traffic models
address phenomena that a static allocation cannot fully describe
\parencite{wang2018}.

Behavioural evidence also matters. The truck-driver study of \textcite{feng2013}
shows that route preferences can vary with vehicle and context.
\Textcite{tahlyan2017} evaluate alternative route sets and route-choice
predictions against truck GPS observations. The choice of plausible alternatives,
including the avoidance of irrelevant routes, affects predictive performance.
Adding random variation to path costs does not by itself establish a calibrated
representation of those preferences.

\Textcite{sakai2020} provide a contrasting integrated configuration:
SimMobility Freight connects commodity contracts, daily logistics, route choice
and parking to a shared passenger--freight traffic simulation. Its illustrative
policies can change both freight and passenger movements. Model components use
observations from several cities, some assumed parameters and local Singapore
count calibration. This makes traffic feedback available for different questions
from supplied-cost loading, while leaving the empirical support for each
component and its transfer to be assessed separately.

The state of the art therefore offers a range of approaches rather than a single
required assignment method. Selection depends on the importance of congestion,
route diversity and temporal effects, alongside the evidence available to model
them. Counts, travel times and observed paths provide complementary validation
targets. Agreement in total traffic or successful numerical convergence cannot
alone establish that the distribution of vehicles across the network is
realistic.

\paragraph{Data requirements and feasibility.}
A usable road graph and suitable travel costs are prerequisites for loading;
truck restrictions, vehicle classes, speeds and observation definitions affect
interpretation. Congestion feedback additionally requires a credible account
of other traffic and network performance. Counts used in calibration should
not also be presented as independent validation. In SimMobility, observations
from several cities, assumed parameters and local count calibration provide
different levels of support for the integrated formulation
\parencite{sakai2020}.

Energy- or slope-informed path objectives need compatible operating inputs.
The Swedish application of \textcite{dias2026} combines operator records,
network and elevation information with slope-aware routing. This demonstrates
a documented operational configuration; it does not establish that equivalent
elevation, fleet and logistics evidence is generally available or that an
energy-efficient path predicts observed route choice.


\paragraph{MASS-GT reference.}
The published network component performs all-or-nothing loading using
least-cost routes on a supplied congested network \parencite{massgt2025}.
Congested travel conditions therefore enter from a transport-model configuration;
this is not evidence of an endogenous freight--background-traffic equilibrium.
The paper assesses network flows using loop detectors and camera counts and
reports limitations in the comparability of these observations. Network-level
agreement should consequently not be extended to every route or logistics
mechanism.

\subsubsection{Environmental assessment}\label{sec:f7}
Inventory approaches combine vehicle activity with emission factors appropriate
to the fleet and operating conditions. The EMEP/EEA guidance provides an
established basis for this accounting and distinguishes hot exhaust,
cold-start effects and non-exhaust sources
\parencite{eeaRoad2026,eeaWear2026}. The consulted sources are updated chapters of
the 2023 guidebook. Their applicability depends on vehicle technology, activity
and the processes covered; tyre and brake wear, for example, cannot be inferred
from an exhaust-only estimate.

The literature offers different ways to describe operating conditions.
Average-speed approaches use representative driving conditions, traffic-situation
approaches distinguish road and service conditions, and more detailed models
represent operating modes, power demand or vehicle trajectories. Physical,
data-driven and hybrid energy models also differ in required inputs and intended
use \parencite{hua2025}. Greater detail is useful when it captures a relevant
phenomenon and can be supported by credible data, rather than simply because it
produces a more elaborate calculation.

Comparative studies demonstrate the importance of the whole modelling chain.
\Textcite{tirico2025} compare traffic and emission models at several spatial
scales. Because differences in traffic representation and assignment are
retained, their results cannot be attributed entirely to the choice of emission
factors; experimental validation of the chain remains future work. In a
heavy-truck fuel study, \textcite{wang2026} assess prediction on vehicles excluded
from fitting. Their outcome is reconstructed from on-board diagnostics, and some
predictors would not necessarily be available before a scenario is simulated.
These are different kinds of evidence from direct pollutant measurements.

Measured-validation evidence clarifies the interpretation of such comparisons.
\Textcite{smit2010} synthesize fifty historical studies using laboratory,
on-board, tunnel, remote-sensing, ambient-concentration and mass-balance
methods. They find no evidence in that sample that greater model complexity
systematically improves prediction errors. Validation designs observe different
processes and scales, and concentration-based assessments also depend on
dispersion and background assumptions. These pre-2010 findings motivate matching
the validation design to the claim; they do not rank current model versions or
provide error bounds for a new implementation.

\Textcite{ligterink2012} demonstrate an empirical payload-sensitive formulation
using on-road measurements from seven Euro-V trucks. Speed, gross mass and
specific power connect vehicle operation to emissions. The work supports
explicit consideration of loading and technology, but its small historical fleet
and fitted driving data limit transfer. It does not establish the accuracy of
current fleets or a universal grade adjustment.

The Swedish case of \textcite{dias2026} illustrates the coupling of logistics,
routing and emissions at an operational scale. It combines operator shipment
data with road-network and elevation information, slope-aware routing and
COPERT-based factors, and considers exhaust and non-exhaust particulate matter.
Its local fleet and logistics configurations delimit interpretation: it is an
application example, not evidence that comparable operator data are generally
available or that its scenario effects transfer unchanged to other cities.

The established requirement is consistency between activity, fleet, operating
conditions and the environmental outcome being estimated. Fuel consumption,
greenhouse gases and regulated pollutants are related but distinct targets.
Non-exhaust emissions, upstream energy production and atmospheric exposure add
further boundaries. A suitable assessment must state which processes it covers
and validate them with corresponding observations, rather than interpreting
model differences as proof that one approach is more accurate.

\paragraph{Data requirements and feasibility.}
Factor applicability depends on the fleet technology, operating conditions,
pollutants and processes represented. Payload, speed and gradient inputs require
compatible units and resolution; fleet assumptions remain assumptions even
when official coefficients are used. The Swedish case relies on
carrier-provided records \parencite{dias2026}. Its data access and local
operating configuration do not establish general feasibility in other areas.
Measured fuel consumption, exhaust tests and concentration observations also
support different validation claims \parencite{smit2010}. Factor substitution
or inter-model agreement cannot replace measurements appropriate to the outcome.


\paragraph{MASS-GT reference.}
The shipment-based method of \textcite{thoen2020emissions} associates vehicle
and link emissions with transported shipments and compares road/load
segmentation with coarser averages. Its results show sensitivity of spatial
estimates to aggregation, not validation against measured pollutants.
The 2025 account uses STREAM-based emission factors by vehicle category,
road category and loading rate, interpolating between empty and full loads
\parencite{massgt2025}. Its traffic-count assessment does not independently
validate pollutant totals. The earlier Transport Policy application explicitly
identified emissions calculation as further work \parencite{massgt2022}.
Thus the emissions anchor must specify its version and factors; changing those
factors later establishes a methodological difference, whose accuracy requires
separate evidence.

\subsection{Integrated findings and evidence limits}\label{sec:discussion}
Across the seven functional responsibilities, the strongest common ground concerns the objects
being modelled and the consistency required between them: goods and vehicle
visits are distinct; logistics choices connect demand to activity; vehicle
movements must be feasible; and environmental estimates depend on the fleet,
operating conditions and processes covered. These principles constrain a model
without prescribing one estimator, routing heuristic or assignment method.

The inspected evidence supports conditional alternatives rather than a universal
ladder of sophistication. A joint choice structure is not always the best-fitting
one; an optimized route need not predict observed behaviour; and detailed
emissions calculations cannot compensate for unsupported fleet or traffic
inputs. Data feasibility is therefore part of methodological suitability, not
only an implementation concern. The appropriate comparison depends on the
question, available information and validation target.

\paragraph{Feasibility of the combined model.}
The separate requirements must also be feasible together. The FHWA inventory
and handbook distinguish data used to run, estimate, calibrate and validate
models at different resolutions \parencite{fhwa2018,fhwa2017}. The UK FAME
peer review likewise distinguishes inputs assumed by a framework from data
available in adequate form \parencite{williams2024}. Compatibility of
classifications, reference years and populations matters across responsibilities,
not only inside each formulation.

The reviewed cases do not establish one minimum data package available in most
study areas. Public aggregates, locally obtained observations and explicitly
estimated or synthetic information support conditional configurations.
Borrowed coefficients, proxies and finer synthetic detail need justification
and uncertainty assessment; they do not automatically replace independent
calibration or validation. Overall feasibility is therefore a joint property
of the selected formulations and their evidence requirements.

MASS-GT makes these dependencies concrete, but its publications document an
evolving model. The fixed implementation reference and the full-system paper
must be reconciled before numerical comparisons. Agreement with that reference
can test implementation correspondence; divergence can identify a changed
assumption or formulation. Neither establishes accuracy without observations
appropriate to the outcome being claimed.

The institutional retrieval strengthens the evidence on aggregate-data touring,
framework transfer, loaded/empty OD consistency, shipment-data enrichment and
measured emissions validation. It also makes the limitations more specific:
calibration can reuse the reported comparison targets; useful spatial detail can
be assumed or synthesized; and scenario benefits depend on imposed operating
and organizational conditions. These distinctions matter more than counting
papers as votes for a single canonical implementation.

Three priority full texts remain unavailable: the data-availability/model-form
book chapter, the urban-logistics/e-commerce data review, and the firm-population
synthesis comparison. Other access requests and incomplete search coverage also
remain explicit. The review remains a documented evidence snapshot, not a
claim of exhaustive field consensus. The functional reorganization changes
presentation and cross-source interpretation, not the underlying retrieval
cohort or the appraisal status of its sources. Positioning NetSy's implemented
formulations and numerical results within this landscape is a subsequent task.
